% !TeX root = ../../Book3.tex
\section{零点计数}
\label{b3:sec:0304}

函数值绕原点转了多少圈，可以通过 \(f'/f\) 的积分读出。零点和极点分别使辐角增加和减少；小扰动若不能把边界像推过原点，内部零点总数就不会变化。

\subsection{辐角原理}
\label{b3:subsec:030401}

对不经过原点的闭连续路径 \(\eta\)，可沿参数区间选连续辐角 \(\theta\)。具体地，把区间细分，使每段像位于一个不包含0的对数圆盘，在相邻端点给辐角加适当的 \(2\pi\) 整数以接合。因首尾函数值相同，\(\theta(b)-\theta(a)\) 必为 \(2\pi\) 的整数倍，其商称为\term[raoshu]{绕数}[winding number]。若路径分段光滑，局部对数求导后相加，得到
\[
 \operatorname{Ind}(\eta,0)=\frac1{2\pi\mathrm{i}}\int_\eta\frac{\mathrm{d}w}{w}.
\]
\term[fujiao-yuanli]{辐角原理}[argument principle]把该整数转成零点计数。
\begin{theorem}[辐角原理]\label{b3:thm:03-argument-principle}
设 \(f\) 在有界分段光滑 Jordan 区域 \(V\) 的闭包邻域亚纯，边界无零点也无极点。则
\[
 \frac1{2\pi\mathrm{i}}\int_{\partial V}\frac{f'(z)}{f(z)}\,\mathrm{d}z=N-P,
\]
其中 \(N\) 与 \(P\) 分别是内部零点和极点的总重数。
\end{theorem}
\begin{proof}
紧性与边界无零极点使内部只有有限个零极点。在零点 \(a\) 的 \(m\) 阶分解中，\(f=(z-a)^mg\)，故 \(f'/f=m/(z-a)+g'/g\)，留数为 \(m\)。在 \(m\) 阶极点处同样得到留数 \(-m\)。留数定理相加给出公式；链式法则说明左端也等于像曲线 \(f\circ\partial V\) 绕0的绕数。
\end{proof}

\subsection{Rouché 定理}
\label{b3:subsec:030402}

\term[rouche-dingli]{Rouché 定理}[Rouché's theorem]允许把难计数的函数与简单函数比较。
\begin{theorem}[Rouché]\label{b3:thm:03-rouche}
设 \(f,g\) 在有界分段光滑 Jordan 区域 \(V\) 的闭包邻域全纯。若边界处 \(|g|<|f|\)，则 \(f\) 与 \(f+g\) 在 \(V\) 内的零点总重数相同。
\end{theorem}
\begin{proof}
对 \(0\le t\le1\)，令 \(f_t=f+tg\)。边界上 \(|f_t|\ge|f|-|g|>0\)，紧边界使右端有共同正下界。辐角原理给出的数
\(N(t)=(2\pi\mathrm{i})^{-1}\int_{\partial V}f_t'/f_t\,\mathrm{d}z\)
依 \(t\) 连续，又总为整数，故在区间上为常数。取 \(t=0,1\) 即得结论。
\end{proof}
例如在 \(|z|=2\) 上，\(|z^5|=32>2+1\)，故 \(z^5+z+1\) 在盘内有5个零点。严格不等式排除了整个变形中边界出现零点。若只有 \(|g|\le|f|\)，边界可能产生零点，开域中的计数就可能改变。

\ProseParagraphBreak

Rouché 的比较方向由边界决定，不能只看多项式的次数。在 \(|z|=1\) 上，\(z^4+3z+1\) 应比较 \(3z\)，因为余项模不超过2；在 \(|z|=2\) 上，\(z^4\) 的模为16，大于 \(3|z|+1=7\)，应比较最高次项。因此同一个多项式在单位盘有1个零点，在半径2盘有4个零点，其余3个位于两圆之间，均计重数。不同半径的零点计数可以相减定位环域中的零点，但边界必须各自无零点。

\ProseParagraphBreak

参数变化时，若一个紧参数区间内的边界函数始终无零点，零点数就保持不变。共同正下界来自参数区间与边界乘积的紧性。如果参数范围不紧，最低模可能趋于0，即使每个固定参数都无边界零点，也不能得到统一的扰动余量。实际使用时可局部按参数取一个小区间，再连通地传播整数计数；不必先找到一个覆盖所有无界参数的误差常数。

\subsection{重数}
\label{b3:subsec:030403}

零点个数若不计重数，就不具备稳定性。函数 \(z^2-\varepsilon\) 在 \(\varepsilon\ne0\) 时有两个不同零点，极限 \(z^2\) 只有一个零点，然而总重数始终为2。
\begin{proposition}[零点局部稳定]\label{b3:prop:03-local-zero-stability}
设 \(f\) 在 \(a\) 处有 \(m\) 阶零点。在充分小的闭圆盘 \(\overline{D(a,r)}\) 上，若全纯函数 \(h\) 满足边界估计 \(|h-f|<\min_{|z-a|=r}|f(z)|\)，则 \(h\) 在圆盘内恰有 \(m\) 个零点，计重数。
\end{proposition}
\begin{proof}
由孤立零点定理选圆盘，使 \(f\) 在其中只有中心零点且边界无零点。给定边界估计就是 Rouché 定理对 \(f\) 与 \(h-f\) 的条件，故两者零点总重数相同。
\end{proof}
若希望定位单个零点，先在每个已知零点周围选互不相交的小圆盘，再分别估计。只在一个大圆上计总数，不会自动给出各个零点的精确位置。

Rouché 定理中的严格不等式只需在边界成立，不要求内部成立。设 \(p(z)=z^5+10z+1\)，在单位圆上以 \(10z\) 为主项，因为 \(|z^5+1|\le2<10\)，可知内部有一个零点；在半径2的圆上，\(|z^5|=32>21\ge|10z+1|\)，内部有五个零点。因此其余四个零点落在 \(1<|z|<2\)。这比求根公式更直接，并且适用于无法写出根的高次多项式。

\ProseParagraphBreak
若把严格不等式放宽为处处等号允许，结论会失败。单位圆上 \(|z|=|1|\)，但 \(z\) 在内部有一个零点，常数1没有零点。严格性保障整个变形 \(f+t g\) 在边界无零；一旦某个参数的零点经过边界，整数零点数便可能跳变。实际计算通常先在紧边界上寻找一个正的间隙，再允许小参数扰动。这个间隙是解析问题中控制根的稳定性的数值。

\subsection{对数导数}
\label{b3:subsec:030404}

对数导数还可以提取带权的零点信息。
\begin{proposition}\label{b3:prop:03-weighted-zeros}
在辐角原理的条件下，若 \(h\) 在 \(\overline V\) 的邻域全纯，则
\[
 \frac1{2\pi\mathrm{i}}\int_{\partial V}h(z)\cdot\frac{f'(z)}{f(z)}\,\mathrm{d}z
 =\sum_a m_a h(a)-\sum_b n_b h(b),
\]
两和分别遍历内部零点与极点，\(m_a,n_b\) 为阶数。
\end{proposition}
\begin{proof}
对数导数在 \(a\) 的主部只有 \(m_a/(z-a)\)，乘全纯 \(h\) 后留数为 \(m_ah(a)\)；极点处相同计算带负号。直接应用留数定理。
\end{proof}
取 \(h(z)=z\) 得到零点之和与极点之和的差；取 \(h(z)=z^k\) 得到幂和。若 \(f\) 在整个区域无零点，\(f'/f\) 本身全纯；它有没有全局原函数，还须看所有闭曲线周期是否为零。\chapref{b3:ch:05}将用这一条件构造全纯对数。

局部零点稳定性始终按重数计数。不同简单根可以合成重根，重根也可分裂；只有在分别选定无零边界的小圆盘之后，才有可持续追踪的局部零点总数。
